\documentclass[10pt]{article}
\usepackage[T1]{fontenc}
\usepackage[utf8]{inputenc}
\usepackage[french]{babel}
\usepackage{amsmath,amssymb,mathrsfs,mathpazo}
\usepackage{geometry}
\geometry{a4paper,margin=2cm}
\usepackage{graphicx,array,tabularx,tabularray,colortbl,multirow,multicol,booktabs,enumitem,pifont,ulem,fancybox,diagbox,cancel}
\usepackage{pstricks,pstricks-add,pst-plot,pst-tree,pst-node,pst-text,pst-func,pst-math,pst-bspline,pst-3dplot,pst-eucl,pst-solides3d}
\let\pscalculate\undefined
\usepackage{pst-bezier}
\usepackage{tikz,pgfplots,tkz-tab}
\usetikzlibrary{arrows,arrows.meta,patterns,calc,positioning,trees,decorations.markings,shapes.geometric,shapes.misc,shapes.arrows,intersections,decorations.pathmorphing,decorations.pathreplacing,backgrounds,shadows,babel}
\pgfplotsset{compat=1.18}
\usepgfplotslibrary{fillbetween}
\usepackage[np]{numprint}
\usepackage{listings,makecell,pas-tableur,twoopt}
\let\barre\undefined
\usepackage[pstricks]{tablvar}
\definecolor{darkgreen}{RGB}{0,100,0}
\providecommand{\red}{\color{red}}
\providecommand{\blue}{\color{blue}}
\providecommand{\green}{\color{green}}
\providecommand{\black}{\color{black}}
\providecommand{\white}{\color{white}}
\providecommand{\gray}{\color{gray}}
\providecommand{\cyan}{\color{cyan}}
\providecommand{\magenta}{\color{magenta}}
\providecommand{\yellow}{\color{yellow}}
\providecommand{\violet}{\color{violet}}
\providecommand{\orange}{\color{orange}}
\providecommand{\euro}{\texteuro}
\providecommand{\EUR}{\texteuro}
\providecommand{\ieme}{\textsuperscript{e}}
\providecommand{\degres}{\ensuremath{^\circ}}
\providecommand{\degree}{\ensuremath{^\circ}}
\providecommand{\vv}[1]{\overrightarrow{#1}}
\providecommand{\Coord}[1]{#1}
\providecommand{\decosix}{}
\providecommand{\starredbullet}{\ensuremath{\bullet}}
\pagestyle{empty}
\setlength{\parindent}{0pt}
\usepackage{adjustbox}\DeclareUnicodeCharacter{2236}{:}\DeclareUnicodeCharacter{2190}{\ensuremath{\leftarrow}}
\providecommand{\R}{\mathbb{R}}
\providecommand{\N}{\mathbb{N}}
\providecommand{\Z}{\mathbb{Z}}
\providecommand{\Q}{\mathbb{Q}}
\providecommand{\D}{\mathbb{D}}
\providecommand{\C}{\mathbb{C}}
\providecommand{\e}{{\mathrm{e}}}
\providecommand{\vect}[1]{\overrightarrow{#1}}
\providecommand{\vectt}[1]{\overrightarrow{\mathrm{#1}}}
\providecommand{\vvt}[1]{\overrightarrow{\mathrm{#1}}}
\providecommand{\np}[2][]{{#2~#1}}
\providecommand{\vv}[1]{{\overrightarrow{#1}}}
\renewcommand{\d}{{\mathrm{d}}}
\providecommand{\degres}{{^\circ}}
\providecommand{\degre}{{^\circ}}
\providecommand{\blue}{\color{blue}}
\providecommand{\red}{\color{red}}
\providecommand{\enskip}{\quad}
\providecommand{\llbracket}{\lbrack\!\lbrack}
\providecommand{\rrbracket}{\rbrack\!\rbrack}
\providecommand{\euro}{{\text{€}}}
\providecommand{\no}{{\text{n°}}}
\providecommand{\interval}[2]{{[#1;#2]}}
\providecommand{\intervFO}[2]{{[#1;#2[}}
\providecommand{\dotfill}{\ldots}

\providecommand{\ds}{\displaystyle}
\providecommand{\textding}[1]{#1}
\providecommand{\Oij}{\(\mathrm{O};\vec{\imath},\vec{\jmath}\)}
\providecommand{\Oijk}{\(\mathrm{O};\vec{\imath},\vec{\jmath},\vec{k}\)}
\providecommand{\Ouv}{\(\mathrm{O};\vec{u},\vec{v}\)}
\providecommand{\intervFF}[2]{\ensuremath{[#1;#2]}}
\renewcommand{\euro}{\ifmmode\text{\texteuro}\else\texteuro\fi}
\DeclareUnicodeCharacter{00B0}{\ifmmode^{\circ}\else\textdegree\fi}
\DeclareUnicodeCharacter{20AC}{\ifmmode\text{\texteuro}\else\texteuro\fi}
% Énoncés : Éducation nationale. Transcription et corrigés : APMEP, Denis Vergès et contributeurs. Découpage et classement : CoopMaths, Sébastien Lozano, Éric Elter et contributeurs. Adaptation de la présentation : Maths Explorer.
% CC BY-SA : https://coopmaths.fr/www/about/
% https://coopmaths.fr/alea/static/bac/2023/tex/bac_2023_03_sujet1_etrangers_groupe1_4_cor.tex
\begin{document}
\noindent\textit{APMEP / CoopMaths. CC BY-SA. Présentation adaptée par Maths Explorer.}\par\medskip
\bigskip

\textbf{Affirmation 1 : Fausse}

Étudions les variations de la fonction $f$.

$f$ est dérivable sur $ [0~;~+\infty[ $, en tant que composée de fonctions dérivables sur $\R$.

Pour tout $t$ réel positif, on pose : $u(t) = -0,5t^2 + t+2$, donc on a $u'(t) = -t + 1$.

On a donc, pour tout réel $t$ positif : $f'(t) = 0 - (-t + 1)\e^{-0,5t^2 + t + 2} = (t-1)\e^{-0,5t^2 + t + 2} $

La fonction exp étant à valeurs strictement positives sur $\R$, le signe de $f'(t)$ est le signe de $(t-1)$ : c'est donc strictement négatif sur l'intervalle $ [0~;~1[ $ et strictement positif sur $ ]1~;~+\infty[ $.

$f$ est donc strictement décroissante sur l'intervalle $[ 0~;~1]$ : c'est pourquoi il est faux de dire que la population augmente en permanence : au cours de la première heure, la population décroit, strictement.

\medskip
\textbf{Affirmation 2 : Fausse}

La fonction exponentielle est à valeurs strictement positives, donc :

$\begin{aligned} \forall t \in \R^+, \quad \e^{-0,5t^2 + t + 2} > 0 &\implies -\e^{-0,5t^2 + t + 2}<0\\
&\implies \e^3-\e^{-0,5t^2 + t + 2}<\e^3\\
&\implies f(t) < \e^3\\ \end{aligned}$

Or $\e^3 \approx 20,086$ (à $ 10^{-3} $ près, arrondi par excès), donc la fonction $f$ est majorée par 20,086, ce qui s'interprète en disant que la population restera inférieure à 20,086 milliers d'individus, soit \np{20086} bactéries : la population ne dépassera donc jamais \np{21000}.

\medskip

\textbf{Affirmation 3 : Vraie}

Une population de bactéries de \np{10000} individus à l'instant $t$, c'est équivalent à avoir $f(t) = 10$.

On cherche donc à savoir combien d'antécédents $10$ a par la fonction $f$.

En reprenant l'étude des variations de $f$ : on a établi que $f$ était strictement décroissante sur $[0~;~1]$, puis strictement croissante sur $ [1~;~+\infty[ $.

Étudions la limite de $f$ en $+\infty$.

$\lim\limits_{t \to +\infty} -0,5t^2 + t +2 = \lim\limits_{t \to +\infty} -0,5t^2 \left(1 - \dfrac{2}{t} - \dfrac{4}{t}\right) = -\infty$, par limite de quotient, somme et produit de suites de références.

Par composition, en posant $y = -0,5t^2 + t +2$, on a $\lim\limits_{t \to+\infty} \e^{-0,5t^2 + t +2} = \lim\limits_{y\to -\infty} \e^y = 0$

Finalement, par limite de la somme $\lim\limits_{t \to +\infty} f(t) = \e^3$.

Pour rappel, on a $\e^3 \approx 20,086 > 10$.

Par ailleurs : $f(0) = \e^3 - \e^2 \approx 12,696 >10$ et $f(1) = \e^3 - \e^{2,5} \approx 7,903 < 10$.

On peut donc dire :

\begin{itemize}
\item Sur $ [0~;~1] $, $f$ est une fonction continue (car dérivable), et strictement décroissante, et 10 est une valeur intermédiaire entre $f(0)$ et $f(1)$, donc en vertu du corollaire au théorème des valeurs intermédiaires appliqué aux fonctions strictement monotones, l'équation $f(t)=10$ admet une unique solution sur l'intervalle $ [0~;~1] $.
\item Sur $ ]1~;~+\infty[ $, $f$ est une fonction continue (car dérivable), et strictement croissante, et 10 est une valeur intermédiaire strictement comprise entre $\lim\limits_{t \to 1}f(t)=f(1)$ et $\lim\limits_{t\to +\infty}f(t)$, donc en vertu du corollaire au théorème des valeurs intermédiaires appliqué aux fonctions strictement monotones, l'équation $f(t)=10$ admet une unique solution sur l'intervalle $]1~;~+\infty[$.
\end{itemize}

Finalement, l'équation $f(t) = 10$ admet exactement deux solutions sur $\R^+$ et donc il y a exactement deux instants où la population de bactéries aura un effectif de \np{10000} individus.

\end{document}
